\documentclass[11pt,a4paper]{article}
% =====================================================================
%  Gemeinsame Vorlage der Arbeitsblätter
%  "Codierung" – Informatik 11 (NTG)
%  (Antwortplatz-Makros wie informatik/10/java/arbeitsblaetter/ab-vorlage.tex)
%  Einbinden in jedem Blatt direkt nach \documentclass: % =====================================================================
%  Gemeinsame Vorlage der Arbeitsblätter
%  "Codierung" – Informatik 11 (NTG)
%  (Antwortplatz-Makros wie informatik/10/java/arbeitsblaetter/ab-vorlage.tex)
%  Einbinden in jedem Blatt direkt nach \documentclass: % =====================================================================
%  Gemeinsame Vorlage der Arbeitsblätter
%  "Codierung" – Informatik 11 (NTG)
%  (Antwortplatz-Makros wie informatik/10/java/arbeitsblaetter/ab-vorlage.tex)
%  Einbinden in jedem Blatt direkt nach \documentclass: \input{ab-vorlage}
%  Übersetzen mit XeLaTeX, LuaLaTeX, pdfLaTeX oder Tectonic.
%
%  Blatt:   tectonic ab-4-1-zahlensysteme.tex   (oder xelatex)
%  Lösung:  tectonic ab-4-1-zahlensysteme-lsg.tex
%           (Datei enthält nur: \def\mitloesung{}\input{ab-4-1-zahlensysteme};
%            füllt alle Lücken rot aus; siehe \lsg, \lsgrest, \lsglinien,
%            \lsgfeld, \lsgkariert, \lsgkarobox, \lsgtext, \lsgoder)
%  Lösungseinträge belegen nie mehr Platz als die Lücke im Blatt.
% =====================================================================
\usepackage[a4paper,left=18mm,right=18mm,top=26mm,bottom=17mm,headheight=15mm,headsep=4mm,footskip=8mm]{geometry}
\usepackage{iftex}
\ifPDFTeX
  \usepackage[T1]{fontenc}
  \usepackage[utf8]{inputenc}
  \usepackage{helvet}
  \usepackage[scaled=0.86]{beramono}
  \renewcommand{\familydefault}{\sfdefault}
\else
  \usepackage{fontspec}
  \setmainfont{texgyreheros}[Extension=.otf,UprightFont=*-regular,BoldFont=*-bold,ItalicFont=*-italic,BoldItalicFont=*-bolditalic]
  \setmonofont{DejaVuSansMono}[Extension=.ttf,UprightFont=*,BoldFont=*-Bold,ItalicFont=*-Oblique,BoldItalicFont=*-BoldOblique,Scale=0.86]
\fi
\usepackage[ngerman,shorthands=off]{babel}
\usepackage{xcolor,graphicx,tabularx,array,enumitem,fancyhdr,tikz,listings,multicol,calc,etoolbox,ragged2e,colortbl,xspace,pifont}
\usepackage[most]{tcolorbox}
\usetikzlibrary{arrows.meta,positioning,calc,decorations.pathreplacing,shapes.geometric,shapes.symbols,fit}
\usepackage[hidelinks]{hyperref}
\setlength{\parindent}{0pt}
\setlength{\parskip}{3pt}
\setlist{nosep,leftmargin=*}

% ---------- Lösungsschalter ----------
\newif\ifloesung
\ifdefined\mitloesung\loesungtrue\fi
\newcommand{\lsgfarbe}{\color{red!75!black}}

% ---------- Kopf- und Fußzeile (einheitliche Form) ----------
\newcommand{\lehrkraft}{Hau}
\newcommand{\themenbereich}{4. Codierung}
\newcommand{\abnummer}{}
\pagestyle{fancy}\fancyhf{}
\renewcommand{\headrulewidth}{0.4pt}
\fancyhead[L]{\small Klasse 11\\Datum:}
\fancyhead[C]{\small Themenbereich:\\\themenbereich}
\fancyhead[R]{\small \lehrkraft\ifloesung\\{\lsgfarbe\bfseries Lösung}\fi}
\fancyfoot[C]{\scriptsize edu-mrh.de\,·\,Informatik 11\,·\,Blatt \abnummer\ifloesung\ (Lösung)\fi\,·\,Seite \thepage}
\newcommand{\abtitel}[2]{\renewcommand{\abnummer}{#1}\begin{center}\Large\bfseries #1\quad\underline{#2}\end{center}\vspace{-1mm}}

% ---------- Boxen ----------
\newenvironment{aufgabe}[2][]{\begin{tcolorbox}[enhanced,breakable,colback=white,colframe=black,boxrule=0.7pt,arc=3mm,left=2mm,right=2mm,top=1.5mm,bottom=1.5mm,before skip=2.5mm,after skip=2.5mm]\textbf{Aufgabe #2}\ifstrempty{#1}{}{ \textit{#1}}\par\vspace{0.6mm}}{\end{tcolorbox}}
\newtcolorbox{merke}[1][Merke]{enhanced,breakable,colback=green!5,colframe=green!40!black,boxrule=0.8pt,arc=2mm,left=2mm,right=2mm,top=1.5mm,bottom=1.5mm,before skip=2.5mm,after skip=2.5mm,title={#1},fonttitle=\bfseries,coltitle=white}
\newtcolorbox{hinweis}{enhanced,breakable,colback=yellow!12,colframe=orange!70!black,boxrule=0.6pt,arc=2mm,left=2mm,right=2mm,top=1mm,bottom=1mm,before skip=2mm,after skip=2mm}
\newcommand{\web}[1]{{\small\color{gray}\textit{Website: #1}}}
% Kasten für optionale Aufgaben
\newtcolorbox{optional}[1][Optional]{enhanced,breakable,colback=blue!3,colframe=blue!45!black,boxrule=0.6pt,arc=2mm,left=2mm,right=2mm,top=1.5mm,bottom=1.5mm,before skip=2.5mm,after skip=2.5mm,title={#1},fonttitle=\bfseries\small,coltitle=white}
% Zahl mit Basis: \zs{1011}{2}
\newcommand{\zs}[2]{#1\textsubscript{#2}}

% ---------- Lücken, Linien, Karos ----------
% Einheitliche Maße für alle Blätter der Sequenz:
%  \linien{n}       n Schreiblinien, Abstand 8 mm, erste Linie 8 mm unter dem Text
%  \kariert{n}      Karoraster 5 mm, n Kästchen hoch, nur ganze Kästchen, mittig
%  \karobox{b}{h}   Karoraster b x h, auf ganze Kästchen gerundet
%  \luecke[l]       Lücke im Fließtext (Standard 4,5 cm, Fachbegriffe mind. 3,5 cm)
%  (jeweils mit Lösungsvariante \lsg…, siehe Kopf der Datei)
%  \restzeile       Lücke bis zum Zeilenende
%  \lueckentext     am Anfang eines Absatzes/Kastens mit Lücken: größerer
%                   Zeilenabstand und Flattersatz, damit man in die Lücken schreiben kann
%  \linienfeld{n}   umrahmtes Feld mit n Schreiblinien (z. B. für Java-Code)
%  \schreibzeile    Strut für Tabellenzellen, in die geschrieben wird (8 mm hoch)
\newlength{\zeilenabstand}\setlength{\zeilenabstand}{8mm}
\newlength{\karo}\setlength{\karo}{5mm}
\tikzset{karo/.style={draw=gray!50,line width=0.3pt},schreiblinie/.style={draw=black!65,line width=0.35pt}}
\newlength{\lsgbreite}
% Lücke: \luecke[Breite] / mit Lösung \lsg[Breite]{Text} (zu langer Text wird verkleinert)
\newcommand{\lsg}[2][4.5cm]{\leavevmode\ifloesung\settowidth{\lsgbreite}{\,#2\,}%
  \ifdim\lsgbreite>#1\relax\rlap{\raisebox{0.8pt}{\resizebox{#1}{!}{\lsgfarbe\,#2\,}}}%
  \else\rlap{\raisebox{0.8pt}{\lsgfarbe\,#2}}\fi\fi\rule[-0.6ex]{#1}{0.4pt}\xspace}
\newcommand{\luecke}[1][4.5cm]{\lsg[#1]{}}
% Lücke bis Zeilenende: \restzeile / \lsgrest{Text}
\newcommand{\lsgrest}[1]{\leavevmode\unskip\hspace{1.5mm}\ifloesung\rlap{\raisebox{0.8pt}{\lsgfarbe\,#1}}\fi\leaders\hbox{\rule[-0.6ex]{0.5pt}{0.4pt}}\hfill\null\par}
\newcommand{\restzeile}{\lsgrest{}}
% nur in der Lösung sichtbar (Platz bleibt frei): \lsgtext{…}; Alternativen: \lsgoder{Lösung}{Blatt}
\newcommand{\lsgtext}[1]{\ifloesung\leavevmode{\lsgfarbe#1}\fi}
\newcommand{\lsgoder}[2]{\ifloesung\leavevmode#1\else#2\fi}
\newcommand{\lueckentext}{\linespread{1.5}\selectfont\raggedright}
\newcommand{\schreibzeile}{\rule[-2.8mm]{0pt}{8mm}}
% Schritt-Nummern wie im Code und auf der Website: \nr{1} = ①
\newcommand{\nr}[1]{\ding{\numexpr171+#1\relax}}
% Schreiblinien: \linien{n} / \lsglinien{n}{Text} (Text liegt auf den Linien)
\newcommand{\lsglinien}[2]{\par\vspace{0.5mm}\noindent\begin{tikzpicture}
  \path (0,0) (\linewidth-0.4pt,-#1\zeilenabstand-1.5mm);
  \foreach \i in {1,...,#1}{\draw[schreiblinie] (0.2pt,-\i\zeilenabstand) -- (\linewidth-0.2pt,-\i\zeilenabstand);}
  \ifloesung\node[overlay,anchor=north west,inner sep=0pt] at (1mm,-\zeilenabstand+1pt+5mm) {\parbox[t]{\linewidth-2mm}{\lsgfarbe\raggedright\setlength{\baselineskip}{\zeilenabstand}\rule{0pt}{5mm}#2}};\fi
  \end{tikzpicture}\par}
\newcommand{\linien}[1]{\lsglinien{#1}{}}
% Codefeld mit Linien: \linienfeld{n} / \lsgfeld{n}{Code, Zeilen mit \\ getrennt}
\newcommand{\lsgfeld}[2]{\noindent\begin{tikzpicture}
  \draw[line width=0.4pt] (0.2pt,0) rectangle (\linewidth-0.2pt,-#1\zeilenabstand-3mm);
  \foreach \i in {1,...,#1}{\draw[schreiblinie] (2mm,-\i\zeilenabstand) -- (\linewidth-2mm,-\i\zeilenabstand);}
  \ifloesung\node[overlay,anchor=north west,inner sep=0pt] at (2.5mm,-\zeilenabstand+1pt+5mm) {\parbox[t]{\linewidth-5mm}{\lsgfarbe\ttfamily\small\setlength{\baselineskip}{\zeilenabstand}\rule{0pt}{5mm}#2}};\fi
  \end{tikzpicture}}
\newcommand{\linienfeld}[1]{\lsgfeld{#1}{}}
% Karoraster; Lösungszeichnung als TikZ-Code, Ursprung oben links, y nach unten negativ:
% \kariert{n} / \lsgkariert{n}{TikZ},  \karobox{b}{h} / \lsgkarobox{b}{h}{TikZ}
\newcommand{\karoraster}[3]{\begin{tikzpicture}\draw[karo] (0,0) grid[step=\karo] (#1\karo,#2\karo);
  \ifloesung\begin{scope}[overlay,shift={(0,#2\karo)},red!75!black,text=red!75!black,line width=0.8pt]#3\end{scope}\fi\end{tikzpicture}}
\newcommand{\lsgkariert}[2]{\par\vspace{1.5mm}\noindent\pgfmathtruncatemacro{\karoX}{floor(\linewidth/\karo)}%
  \makebox[\linewidth]{\karoraster{\karoX}{#1}{#2}}\par\vspace{1mm}}
\newcommand{\kariert}[1]{\lsgkariert{#1}{}}
\newcommand{\lsgkarobox}[3]{\pgfmathtruncatemacro{\karoX}{round((#1)/\karo)}\pgfmathtruncatemacro{\karoY}{round((#2)/\karo)}\karoraster{\karoX}{\karoY}{#3}}
\newcommand{\karobox}[2]{\lsgkarobox{#1}{#2}{}}
% Lösungs-Klassenkarte (für Zeichnungen im Karo): \lsgkarte{Name}{Attribute}{Methoden}
\newcommand{\lsgkarte}[3]{{\lsgfarbe\setlength{\arrayrulewidth}{0.6pt}\arrayrulecolor{red!75!black}\begin{tabular}{|l|}\hline\textbf{#1}\\\hline #2\\\hline #3\\\hline\end{tabular}}}
% Lösungs-Objektkarte (in TikZ als Knoteninhalt): \node[objekt] {\lsgobjekt{name : Klasse}{attr = wert\\…}};
\newcommand{\lsgobjekt}[2]{{\lsgfarbe\small\arrayrulecolor{red!75!black}\begin{tabular}{@{}l@{}}\textbf{#1}\\\hline #2\end{tabular}}}
\tikzset{objekt/.style={draw,rounded corners=3mm,inner sep=2mm,anchor=north west}}
\newcommand{\klassenkarte}[3]{\begin{tabular}{|l|}\hline\textbf{#1}\\\hline #2\\\hline #3\\\hline\end{tabular}}
\newcommand{\code}[1]{\texttt{#1}}

% ---------- Java-Listings ----------
\lstset{language=Java,basicstyle=\ttfamily\small,keywordstyle=\bfseries,commentstyle=\color{gray!80!black},stringstyle=\color{black},showstringspaces=false,columns=flexible,keepspaces=true,tabsize=3,frame=single,framerule=0.4pt,rulecolor=\color{black},xleftmargin=2mm,framexleftmargin=1mm,aboveskip=2mm,belowskip=2mm,breaklines=true,
 literate={ä}{{\"a}}1 {ö}{{\"o}}1 {ü}{{\"u}}1 {Ä}{{\"A}}1 {Ö}{{\"O}}1 {Ü}{{\"U}}1 {ß}{{\ss}}1 {–}{{--}}1 {„}{{\glqq}}1 {“}{{\grqq}}1}
\lstnewenvironment{java}[1][]{\lstset{#1}}{}
% Lückencode: größerer Zeilenabstand, Lücken mit (*\luecke[3cm]*) im Code
\lstnewenvironment{javaluecke}[1][]{\lstset{basicstyle=\linespread{1.5}\small\ttfamily,escapeinside={(*}{*)},#1}}{}

%  Übersetzen mit XeLaTeX, LuaLaTeX, pdfLaTeX oder Tectonic.
%
%  Blatt:   tectonic ab-4-1-zahlensysteme.tex   (oder xelatex)
%  Lösung:  tectonic ab-4-1-zahlensysteme-lsg.tex
%           (Datei enthält nur: \def\mitloesung{}\documentclass[11pt,a4paper]{article}
\input{ab-vorlage}
\newcolumntype{C}{>{\centering\arraybackslash}X}
\newcommand{\kopfzelle}[1]{\cellcolor{blue!8}\textbf{#1}}
\begin{document}
\abtitel{4.1}{Zahlensysteme}

\textbf{Das Dezimalsystem}\par
Unser gängiges Zahlensystem zur Basis 10 (griech. \textit{déka}, lat. \textit{decem}: zehn) kennt die zehn Ziffern 0 bis 9. Jede Stelle hat einen festen \textbf{Stellenwert}, eine Potenz von 10.\\[1mm]
Beispiel:\quad $35\,426 = 3\cdot 10\,000 + 5\cdot 1000 + 4\cdot 100 + 2\cdot 10 + 6\cdot 1$\\
\hspace*{2.05cm}$= 3\cdot 10^4 + 5\cdot 10^3 + 4\cdot 10^2 + 2\cdot 10^1 + 6\cdot 10^0$\par\vspace{1mm}
{\renewcommand{\arraystretch}{1.25}
\begin{tabularx}{\linewidth}{|>{\raggedright\arraybackslash}p{2.6cm}|C|C|C|C|C|}\hline
\kopfzelle{Stellenwert} & {\small Zehntausender ($10^4$)} & {\small Tausender ($10^3$)} & {\small Hunderter ($10^2$)} & {\small Zehner ($10^1$)} & {\small Einer ($10^0$)}\\\hline
\kopfzelle{Ziffer} & 3 & 5 & 4 & 2 & 6\\\hline
\end{tabularx}}

\medskip
\textbf{Das Binärsystem (zur Basis \lsg[1.5cm]{2})}\par
{\lueckentext Das Binärsystem kennt nur die Ziffern \lsg[1.5cm]{0} und \lsg[1.5cm]{1}. Eine solche Ziffer heißt \textbf{Bit} (\textit{binary digit}). Wie im Dezimalsystem hat jede Stelle einen festen \textbf{Stellenwert}. Die Stellenwerte werden \textbf{von rechts nach links} größer.\par}
\vspace{1mm}
{\renewcommand{\arraystretch}{1.1}
\begin{tabularx}{\linewidth}{|>{\raggedright\arraybackslash}p{2.6cm}|C|C|C|C|C|}\hline
\kopfzelle{Binärzahl} & \Large 1 & \Large 1 & \Large 1 & \Large 0 & \Large 1\\\hline
\kopfzelle{Stellenwert (Potenz)} & \schreibzeile\lsgtext{$2^4 = 16$} & \lsgtext{$2^3 = 8$} & \lsgtext{$2^2 = 4$} & \lsgtext{$2^1 = 2$} & \lsgtext{$2^0 = 1$}\\\hline
\kopfzelle{Produkt} & \schreibzeile\lsgtext{$1\cdot 16 = 16$} & \lsgtext{$1\cdot 8 = 8$} & \lsgtext{$1\cdot 4 = 4$} & \lsgtext{$0\cdot 2 = 0$} & \lsgtext{$1\cdot 1 = 1$}\\\hline
\kopfzelle{Dezimalzahl} & \multicolumn{5}{c|}{\schreibzeile\lsgtext{$16 + 8 + 4 + 0 + 1 = 29$}}\\\hline
\end{tabularx}}

\begin{merke}[Bit und Byte]
{\lueckentext Mit $n$ Bit lassen sich \lsg[2.5cm]{$2^n$} verschiedene Zahlen darstellen, nämlich 0 bis $2^n - 1$. 8 Bit fasst man zu \textbf{einem Byte} zusammen. Ein Byte kann also \lsg[2cm]{256} verschiedene Werte speichern, von 0 bis \lsg[1.5cm]{255}.\par}
\end{merke}

\begin{minipage}[t]{0.5\linewidth}
\textbf{Das Hexadezimalsystem (zur Basis \lsg[1.5cm]{16})}\par
Hier braucht man 16 Ziffern. Nach der 9 geht es mit Buchstaben weiter:\par\vspace{1.5mm}
{\small\setlength{\tabcolsep}{3pt}\renewcommand{\arraystretch}{1.15}
\begin{tabular}{|l|c|c|c|c|c|c|c|}\hline
\kopfzelle{Ziffer} & 0–9 & A & B & C & D & E & F\\\hline
\kopfzelle{Wert} & 0–9 & \schreibzeile\lsgtext{10} & \lsgtext{11} & \lsgtext{12} & \lsgtext{13} & \lsgtext{14} & \lsgtext{15}\\\hline
\end{tabular}}
\end{minipage}\hfill
\begin{minipage}[t]{0.47\linewidth}
\vspace{0pt}
{\renewcommand{\arraystretch}{1.1}\small
\begin{tabularx}{\linewidth}{|>{\raggedright\arraybackslash}p{2.35cm}|C|C|C|}\hline
\kopfzelle{Hexadezimal\-zahl} & \multicolumn{3}{c|}{\schreibzeile\Large\lsgtext{2\qquad B\qquad C}}\\\hline
\kopfzelle{Stellenwert (Potenz)} & \schreibzeile\lsgtext{$16^2{=}256$} & \lsgtext{$16^1{=}16$} & \lsgtext{$16^0{=}1$}\\\hline
\kopfzelle{Produkt} & \schreibzeile\lsgtext{$2\cdot 256$} & \lsgtext{$11\cdot 16$} & \lsgtext{$12\cdot 1$}\\\hline
\kopfzelle{Dezimalzahl} & \multicolumn{3}{c|}{\schreibzeile\normalsize 700}\\\hline
\end{tabularx}}
\end{minipage}

\medskip
\begin{minipage}[t]{0.45\linewidth}
\textbf{Umrechnung vom Dezimalsystem in ein anderes System} (Zahl $x$)\par\vspace{0.5mm}
\begin{enumerate}[label=\nr{\arabic*},leftmargin=6mm]
\item Finde den größten Stellenwert $S$ des neuen Systems, der nicht größer als $x$ ist.
\item Dividiere $x$ durch $S$ mit Rest $R$. Das ganzzahlige Ergebnis ist die Ziffer an dieser Stelle.
\item Mache mit dem Rest $R$ beim nächstkleineren Stellenwert weiter, bis die Stelle mit $S = 1$ erledigt ist.
\end{enumerate}
\end{minipage}\hfill
\begin{minipage}[t]{0.52\linewidth}
\vspace{0pt}
{\renewcommand{\arraystretch}{1.75}
\begin{tabular}{@{}l@{\ \,}l@{\ \,}l@{}}
\textbf{1. Durchgang:} & $S = 256$ & $700 : 256 = \mathbf{2}$ R $188$\\
\textbf{2. Durchgang:} & $S = $ \lsg[1.2cm]{16} & $188 : $ \lsg[3.2cm]{$16 = 11$ (B) R 12}\\
\textbf{3. Durchgang:} & $S = $ \lsg[1.2cm]{1} & \lsg[3.9cm]{$12 : 1 = 12$ (C) R 0}\\
\textbf{Ergebnis:} & \multicolumn{2}{l}{$700 = $ \lsg[2.5cm]{2BC}\textsubscript{16}}
\end{tabular}}
\end{minipage}

\newpage
\begin{aufgabe}[Bit-Schalter und Trainer]{1}
{\lueckentext\textbf{a)} Stelle mit dem \textbf{Bit-Schalter} auf der Website die Zahlen 5, 18 und 31 ein. Die größte Zahl, die man mit 5 Bit darstellen kann, ist \lsg[1.5cm]{31}, mit 8 Bit \lsg[1.5cm]{255}.\\
\textbf{b)} Bearbeite im \textbf{Zahlensystem-Trainer} (Level „Einfach“) mindestens 10 Aufgaben.\par}
\end{aufgabe}

\begin{aufgabe}[Umrechnen]{2}
Wandle die Zahlen \textbf{1337} und \textbf{4095} ins Hexadezimalsystem \textbf{und} ins Binärsystem um. Rechne mit dem Verfahren von Seite 1.
\lsgkariert{11}{\begin{scope}[font=\small,anchor=west]
\node at (2mm,-2.5mm) {1337 → Hex:\quad $1337 : 256 = 5$ R $57$;\quad $57 : 16 = 3$ R $9$;\quad $9 : 1 = 9$ \quad ⇒\quad 539\textsubscript{16}};
\node at (2mm,-7.5mm) {1337 → Binär: \ $S = 1024$: 1 R 313;\ 512: 0;\ 256: 1 R 57;\ 128: 0;\ 64: 0;\ 32: 1 R 25;};
\node at (2mm,-12.5mm) {\hspace{2.2cm}16: 1 R 9;\ 8: 1 R 1;\ 4: 0;\ 2: 0;\ 1: 1 \quad ⇒\quad 101\,0011\,1001\textsubscript{2}};
\node at (2mm,-22.5mm) {4095 → Hex:\quad $4095 : 256 = 15$ (F) R $255$;\quad $255 : 16 = 15$ (F) R $15$;\quad $15 : 1 = 15$ (F) \quad ⇒\quad FFF\textsubscript{16}};
\node at (2mm,-27.5mm) {4095 → Binär: \ $S = 2048$: 1 R 2047;\ 1024: 1 R 1023;\ 512: 1 R 511;\ \ldots\ jede Stelle 1};
\node at (2mm,-32.5mm) {\hspace{2.2cm}⇒\quad 1111\,1111\,1111\textsubscript{2}\quad (12 Bit; $4095 = 2^{12} - 1$)};
\end{scope}}
{\lueckentext $1337 = $ \lsg[3cm]{539}\textsubscript{16} $=$ \lsg[6cm]{101\,0011\,1001}\textsubscript{2}\\
$4095 = $ \lsg[3cm]{FFF}\textsubscript{16} $=$ \lsg[6cm]{1111\,1111\,1111}\textsubscript{2}\par}
\textit{Zum Nachdenken (mündlich):} Was fällt dir bei 4095 auf? Welche Zahl kommt danach?
\end{aufgabe}

\begin{aufgabe}[Die Abkürzung: Binär $\leftrightarrow$ Hexadezimal]{3}
Eine Hexadezimalziffer steht für genau 4 Bit, denn $16 = 2^4$. Teile eine Binärzahl deshalb \textbf{von rechts} in Vierergruppen und übersetze jede Gruppe einzeln.\\[1mm]
Beispiel:\quad \zs{1101\,0110}{2}\ $\to$\ 1101 = 13 = D,\ \ 0110 = 6\ $\to$\ \zs{D6}{16}\par\vspace{1mm}
{\lueckentext
\textbf{a)} \zs{1111\,0000}{2} = \lsg[1.6cm]{F0}\textsubscript{16}\qquad
\textbf{b)} \zs{1010\,0111}{2} = \lsg[1.6cm]{A7}\textsubscript{16}\qquad
\textbf{c)} \zs{7E}{16} = \lsg[3.2cm]{0111\,1110}\textsubscript{2}\\
\textbf{d)} Kontrolliere so dein Ergebnis aus Aufgabe 2 für 1337.\par}
\end{aufgabe}

\begin{aufgabe}[The Mystery of Cryptocastle]{4}
Öffne das Spiel \textit{The Mystery of Cryptocastle} (Link auf der Website) und bearbeite die Aufgaben in \textbf{„2 Zum alten Burgturm“}. Notiere hier deine Rechnungen und Codes:
\linien{3}
\end{aufgabe}

\vfill
\web{Station 4.1 Zahlensysteme: Bit-Schalter, Umrechnungs-Animation, Trainer}
\end{document}
;
%            füllt alle Lücken rot aus; siehe \lsg, \lsgrest, \lsglinien,
%            \lsgfeld, \lsgkariert, \lsgkarobox, \lsgtext, \lsgoder)
%  Lösungseinträge belegen nie mehr Platz als die Lücke im Blatt.
% =====================================================================
\usepackage[a4paper,left=18mm,right=18mm,top=26mm,bottom=17mm,headheight=15mm,headsep=4mm,footskip=8mm]{geometry}
\usepackage{iftex}
\ifPDFTeX
  \usepackage[T1]{fontenc}
  \usepackage[utf8]{inputenc}
  \usepackage{helvet}
  \usepackage[scaled=0.86]{beramono}
  \renewcommand{\familydefault}{\sfdefault}
\else
  \usepackage{fontspec}
  \setmainfont{texgyreheros}[Extension=.otf,UprightFont=*-regular,BoldFont=*-bold,ItalicFont=*-italic,BoldItalicFont=*-bolditalic]
  \setmonofont{DejaVuSansMono}[Extension=.ttf,UprightFont=*,BoldFont=*-Bold,ItalicFont=*-Oblique,BoldItalicFont=*-BoldOblique,Scale=0.86]
\fi
\usepackage[ngerman,shorthands=off]{babel}
\usepackage{xcolor,graphicx,tabularx,array,enumitem,fancyhdr,tikz,listings,multicol,calc,etoolbox,ragged2e,colortbl,xspace,pifont}
\usepackage[most]{tcolorbox}
\usetikzlibrary{arrows.meta,positioning,calc,decorations.pathreplacing,shapes.geometric,shapes.symbols,fit}
\usepackage[hidelinks]{hyperref}
\setlength{\parindent}{0pt}
\setlength{\parskip}{3pt}
\setlist{nosep,leftmargin=*}

% ---------- Lösungsschalter ----------
\newif\ifloesung
\ifdefined\mitloesung\loesungtrue\fi
\newcommand{\lsgfarbe}{\color{red!75!black}}

% ---------- Kopf- und Fußzeile (einheitliche Form) ----------
\newcommand{\lehrkraft}{Hau}
\newcommand{\themenbereich}{4. Codierung}
\newcommand{\abnummer}{}
\pagestyle{fancy}\fancyhf{}
\renewcommand{\headrulewidth}{0.4pt}
\fancyhead[L]{\small Klasse 11\\Datum:}
\fancyhead[C]{\small Themenbereich:\\\themenbereich}
\fancyhead[R]{\small \lehrkraft\ifloesung\\{\lsgfarbe\bfseries Lösung}\fi}
\fancyfoot[C]{\scriptsize edu-mrh.de\,·\,Informatik 11\,·\,Blatt \abnummer\ifloesung\ (Lösung)\fi\,·\,Seite \thepage}
\newcommand{\abtitel}[2]{\renewcommand{\abnummer}{#1}\begin{center}\Large\bfseries #1\quad\underline{#2}\end{center}\vspace{-1mm}}

% ---------- Boxen ----------
\newenvironment{aufgabe}[2][]{\begin{tcolorbox}[enhanced,breakable,colback=white,colframe=black,boxrule=0.7pt,arc=3mm,left=2mm,right=2mm,top=1.5mm,bottom=1.5mm,before skip=2.5mm,after skip=2.5mm]\textbf{Aufgabe #2}\ifstrempty{#1}{}{ \textit{#1}}\par\vspace{0.6mm}}{\end{tcolorbox}}
\newtcolorbox{merke}[1][Merke]{enhanced,breakable,colback=green!5,colframe=green!40!black,boxrule=0.8pt,arc=2mm,left=2mm,right=2mm,top=1.5mm,bottom=1.5mm,before skip=2.5mm,after skip=2.5mm,title={#1},fonttitle=\bfseries,coltitle=white}
\newtcolorbox{hinweis}{enhanced,breakable,colback=yellow!12,colframe=orange!70!black,boxrule=0.6pt,arc=2mm,left=2mm,right=2mm,top=1mm,bottom=1mm,before skip=2mm,after skip=2mm}
\newcommand{\web}[1]{{\small\color{gray}\textit{Website: #1}}}
% Kasten für optionale Aufgaben
\newtcolorbox{optional}[1][Optional]{enhanced,breakable,colback=blue!3,colframe=blue!45!black,boxrule=0.6pt,arc=2mm,left=2mm,right=2mm,top=1.5mm,bottom=1.5mm,before skip=2.5mm,after skip=2.5mm,title={#1},fonttitle=\bfseries\small,coltitle=white}
% Zahl mit Basis: \zs{1011}{2}
\newcommand{\zs}[2]{#1\textsubscript{#2}}

% ---------- Lücken, Linien, Karos ----------
% Einheitliche Maße für alle Blätter der Sequenz:
%  \linien{n}       n Schreiblinien, Abstand 8 mm, erste Linie 8 mm unter dem Text
%  \kariert{n}      Karoraster 5 mm, n Kästchen hoch, nur ganze Kästchen, mittig
%  \karobox{b}{h}   Karoraster b x h, auf ganze Kästchen gerundet
%  \luecke[l]       Lücke im Fließtext (Standard 4,5 cm, Fachbegriffe mind. 3,5 cm)
%  (jeweils mit Lösungsvariante \lsg…, siehe Kopf der Datei)
%  \restzeile       Lücke bis zum Zeilenende
%  \lueckentext     am Anfang eines Absatzes/Kastens mit Lücken: größerer
%                   Zeilenabstand und Flattersatz, damit man in die Lücken schreiben kann
%  \linienfeld{n}   umrahmtes Feld mit n Schreiblinien (z. B. für Java-Code)
%  \schreibzeile    Strut für Tabellenzellen, in die geschrieben wird (8 mm hoch)
\newlength{\zeilenabstand}\setlength{\zeilenabstand}{8mm}
\newlength{\karo}\setlength{\karo}{5mm}
\tikzset{karo/.style={draw=gray!50,line width=0.3pt},schreiblinie/.style={draw=black!65,line width=0.35pt}}
\newlength{\lsgbreite}
% Lücke: \luecke[Breite] / mit Lösung \lsg[Breite]{Text} (zu langer Text wird verkleinert)
\newcommand{\lsg}[2][4.5cm]{\leavevmode\ifloesung\settowidth{\lsgbreite}{\,#2\,}%
  \ifdim\lsgbreite>#1\relax\rlap{\raisebox{0.8pt}{\resizebox{#1}{!}{\lsgfarbe\,#2\,}}}%
  \else\rlap{\raisebox{0.8pt}{\lsgfarbe\,#2}}\fi\fi\rule[-0.6ex]{#1}{0.4pt}\xspace}
\newcommand{\luecke}[1][4.5cm]{\lsg[#1]{}}
% Lücke bis Zeilenende: \restzeile / \lsgrest{Text}
\newcommand{\lsgrest}[1]{\leavevmode\unskip\hspace{1.5mm}\ifloesung\rlap{\raisebox{0.8pt}{\lsgfarbe\,#1}}\fi\leaders\hbox{\rule[-0.6ex]{0.5pt}{0.4pt}}\hfill\null\par}
\newcommand{\restzeile}{\lsgrest{}}
% nur in der Lösung sichtbar (Platz bleibt frei): \lsgtext{…}; Alternativen: \lsgoder{Lösung}{Blatt}
\newcommand{\lsgtext}[1]{\ifloesung\leavevmode{\lsgfarbe#1}\fi}
\newcommand{\lsgoder}[2]{\ifloesung\leavevmode#1\else#2\fi}
\newcommand{\lueckentext}{\linespread{1.5}\selectfont\raggedright}
\newcommand{\schreibzeile}{\rule[-2.8mm]{0pt}{8mm}}
% Schritt-Nummern wie im Code und auf der Website: \nr{1} = ①
\newcommand{\nr}[1]{\ding{\numexpr171+#1\relax}}
% Schreiblinien: \linien{n} / \lsglinien{n}{Text} (Text liegt auf den Linien)
\newcommand{\lsglinien}[2]{\par\vspace{0.5mm}\noindent\begin{tikzpicture}
  \path (0,0) (\linewidth-0.4pt,-#1\zeilenabstand-1.5mm);
  \foreach \i in {1,...,#1}{\draw[schreiblinie] (0.2pt,-\i\zeilenabstand) -- (\linewidth-0.2pt,-\i\zeilenabstand);}
  \ifloesung\node[overlay,anchor=north west,inner sep=0pt] at (1mm,-\zeilenabstand+1pt+5mm) {\parbox[t]{\linewidth-2mm}{\lsgfarbe\raggedright\setlength{\baselineskip}{\zeilenabstand}\rule{0pt}{5mm}#2}};\fi
  \end{tikzpicture}\par}
\newcommand{\linien}[1]{\lsglinien{#1}{}}
% Codefeld mit Linien: \linienfeld{n} / \lsgfeld{n}{Code, Zeilen mit \\ getrennt}
\newcommand{\lsgfeld}[2]{\noindent\begin{tikzpicture}
  \draw[line width=0.4pt] (0.2pt,0) rectangle (\linewidth-0.2pt,-#1\zeilenabstand-3mm);
  \foreach \i in {1,...,#1}{\draw[schreiblinie] (2mm,-\i\zeilenabstand) -- (\linewidth-2mm,-\i\zeilenabstand);}
  \ifloesung\node[overlay,anchor=north west,inner sep=0pt] at (2.5mm,-\zeilenabstand+1pt+5mm) {\parbox[t]{\linewidth-5mm}{\lsgfarbe\ttfamily\small\setlength{\baselineskip}{\zeilenabstand}\rule{0pt}{5mm}#2}};\fi
  \end{tikzpicture}}
\newcommand{\linienfeld}[1]{\lsgfeld{#1}{}}
% Karoraster; Lösungszeichnung als TikZ-Code, Ursprung oben links, y nach unten negativ:
% \kariert{n} / \lsgkariert{n}{TikZ},  \karobox{b}{h} / \lsgkarobox{b}{h}{TikZ}
\newcommand{\karoraster}[3]{\begin{tikzpicture}\draw[karo] (0,0) grid[step=\karo] (#1\karo,#2\karo);
  \ifloesung\begin{scope}[overlay,shift={(0,#2\karo)},red!75!black,text=red!75!black,line width=0.8pt]#3\end{scope}\fi\end{tikzpicture}}
\newcommand{\lsgkariert}[2]{\par\vspace{1.5mm}\noindent\pgfmathtruncatemacro{\karoX}{floor(\linewidth/\karo)}%
  \makebox[\linewidth]{\karoraster{\karoX}{#1}{#2}}\par\vspace{1mm}}
\newcommand{\kariert}[1]{\lsgkariert{#1}{}}
\newcommand{\lsgkarobox}[3]{\pgfmathtruncatemacro{\karoX}{round((#1)/\karo)}\pgfmathtruncatemacro{\karoY}{round((#2)/\karo)}\karoraster{\karoX}{\karoY}{#3}}
\newcommand{\karobox}[2]{\lsgkarobox{#1}{#2}{}}
% Lösungs-Klassenkarte (für Zeichnungen im Karo): \lsgkarte{Name}{Attribute}{Methoden}
\newcommand{\lsgkarte}[3]{{\lsgfarbe\setlength{\arrayrulewidth}{0.6pt}\arrayrulecolor{red!75!black}\begin{tabular}{|l|}\hline\textbf{#1}\\\hline #2\\\hline #3\\\hline\end{tabular}}}
% Lösungs-Objektkarte (in TikZ als Knoteninhalt): \node[objekt] {\lsgobjekt{name : Klasse}{attr = wert\\…}};
\newcommand{\lsgobjekt}[2]{{\lsgfarbe\small\arrayrulecolor{red!75!black}\begin{tabular}{@{}l@{}}\textbf{#1}\\\hline #2\end{tabular}}}
\tikzset{objekt/.style={draw,rounded corners=3mm,inner sep=2mm,anchor=north west}}
\newcommand{\klassenkarte}[3]{\begin{tabular}{|l|}\hline\textbf{#1}\\\hline #2\\\hline #3\\\hline\end{tabular}}
\newcommand{\code}[1]{\texttt{#1}}

% ---------- Java-Listings ----------
\lstset{language=Java,basicstyle=\ttfamily\small,keywordstyle=\bfseries,commentstyle=\color{gray!80!black},stringstyle=\color{black},showstringspaces=false,columns=flexible,keepspaces=true,tabsize=3,frame=single,framerule=0.4pt,rulecolor=\color{black},xleftmargin=2mm,framexleftmargin=1mm,aboveskip=2mm,belowskip=2mm,breaklines=true,
 literate={ä}{{\"a}}1 {ö}{{\"o}}1 {ü}{{\"u}}1 {Ä}{{\"A}}1 {Ö}{{\"O}}1 {Ü}{{\"U}}1 {ß}{{\ss}}1 {–}{{--}}1 {„}{{\glqq}}1 {“}{{\grqq}}1}
\lstnewenvironment{java}[1][]{\lstset{#1}}{}
% Lückencode: größerer Zeilenabstand, Lücken mit (*\luecke[3cm]*) im Code
\lstnewenvironment{javaluecke}[1][]{\lstset{basicstyle=\linespread{1.5}\small\ttfamily,escapeinside={(*}{*)},#1}}{}

%  Übersetzen mit XeLaTeX, LuaLaTeX, pdfLaTeX oder Tectonic.
%
%  Blatt:   tectonic ab-4-1-zahlensysteme.tex   (oder xelatex)
%  Lösung:  tectonic ab-4-1-zahlensysteme-lsg.tex
%           (Datei enthält nur: \def\mitloesung{}\documentclass[11pt,a4paper]{article}
% =====================================================================
%  Gemeinsame Vorlage der Arbeitsblätter
%  "Codierung" – Informatik 11 (NTG)
%  (Antwortplatz-Makros wie informatik/10/java/arbeitsblaetter/ab-vorlage.tex)
%  Einbinden in jedem Blatt direkt nach \documentclass: \input{ab-vorlage}
%  Übersetzen mit XeLaTeX, LuaLaTeX, pdfLaTeX oder Tectonic.
%
%  Blatt:   tectonic ab-4-1-zahlensysteme.tex   (oder xelatex)
%  Lösung:  tectonic ab-4-1-zahlensysteme-lsg.tex
%           (Datei enthält nur: \def\mitloesung{}\input{ab-4-1-zahlensysteme};
%            füllt alle Lücken rot aus; siehe \lsg, \lsgrest, \lsglinien,
%            \lsgfeld, \lsgkariert, \lsgkarobox, \lsgtext, \lsgoder)
%  Lösungseinträge belegen nie mehr Platz als die Lücke im Blatt.
% =====================================================================
\usepackage[a4paper,left=18mm,right=18mm,top=26mm,bottom=17mm,headheight=15mm,headsep=4mm,footskip=8mm]{geometry}
\usepackage{iftex}
\ifPDFTeX
  \usepackage[T1]{fontenc}
  \usepackage[utf8]{inputenc}
  \usepackage{helvet}
  \usepackage[scaled=0.86]{beramono}
  \renewcommand{\familydefault}{\sfdefault}
\else
  \usepackage{fontspec}
  \setmainfont{texgyreheros}[Extension=.otf,UprightFont=*-regular,BoldFont=*-bold,ItalicFont=*-italic,BoldItalicFont=*-bolditalic]
  \setmonofont{DejaVuSansMono}[Extension=.ttf,UprightFont=*,BoldFont=*-Bold,ItalicFont=*-Oblique,BoldItalicFont=*-BoldOblique,Scale=0.86]
\fi
\usepackage[ngerman,shorthands=off]{babel}
\usepackage{xcolor,graphicx,tabularx,array,enumitem,fancyhdr,tikz,listings,multicol,calc,etoolbox,ragged2e,colortbl,xspace,pifont}
\usepackage[most]{tcolorbox}
\usetikzlibrary{arrows.meta,positioning,calc,decorations.pathreplacing,shapes.geometric,shapes.symbols,fit}
\usepackage[hidelinks]{hyperref}
\setlength{\parindent}{0pt}
\setlength{\parskip}{3pt}
\setlist{nosep,leftmargin=*}

% ---------- Lösungsschalter ----------
\newif\ifloesung
\ifdefined\mitloesung\loesungtrue\fi
\newcommand{\lsgfarbe}{\color{red!75!black}}

% ---------- Kopf- und Fußzeile (einheitliche Form) ----------
\newcommand{\lehrkraft}{Hau}
\newcommand{\themenbereich}{4. Codierung}
\newcommand{\abnummer}{}
\pagestyle{fancy}\fancyhf{}
\renewcommand{\headrulewidth}{0.4pt}
\fancyhead[L]{\small Klasse 11\\Datum:}
\fancyhead[C]{\small Themenbereich:\\\themenbereich}
\fancyhead[R]{\small \lehrkraft\ifloesung\\{\lsgfarbe\bfseries Lösung}\fi}
\fancyfoot[C]{\scriptsize edu-mrh.de\,·\,Informatik 11\,·\,Blatt \abnummer\ifloesung\ (Lösung)\fi\,·\,Seite \thepage}
\newcommand{\abtitel}[2]{\renewcommand{\abnummer}{#1}\begin{center}\Large\bfseries #1\quad\underline{#2}\end{center}\vspace{-1mm}}

% ---------- Boxen ----------
\newenvironment{aufgabe}[2][]{\begin{tcolorbox}[enhanced,breakable,colback=white,colframe=black,boxrule=0.7pt,arc=3mm,left=2mm,right=2mm,top=1.5mm,bottom=1.5mm,before skip=2.5mm,after skip=2.5mm]\textbf{Aufgabe #2}\ifstrempty{#1}{}{ \textit{#1}}\par\vspace{0.6mm}}{\end{tcolorbox}}
\newtcolorbox{merke}[1][Merke]{enhanced,breakable,colback=green!5,colframe=green!40!black,boxrule=0.8pt,arc=2mm,left=2mm,right=2mm,top=1.5mm,bottom=1.5mm,before skip=2.5mm,after skip=2.5mm,title={#1},fonttitle=\bfseries,coltitle=white}
\newtcolorbox{hinweis}{enhanced,breakable,colback=yellow!12,colframe=orange!70!black,boxrule=0.6pt,arc=2mm,left=2mm,right=2mm,top=1mm,bottom=1mm,before skip=2mm,after skip=2mm}
\newcommand{\web}[1]{{\small\color{gray}\textit{Website: #1}}}
% Kasten für optionale Aufgaben
\newtcolorbox{optional}[1][Optional]{enhanced,breakable,colback=blue!3,colframe=blue!45!black,boxrule=0.6pt,arc=2mm,left=2mm,right=2mm,top=1.5mm,bottom=1.5mm,before skip=2.5mm,after skip=2.5mm,title={#1},fonttitle=\bfseries\small,coltitle=white}
% Zahl mit Basis: \zs{1011}{2}
\newcommand{\zs}[2]{#1\textsubscript{#2}}

% ---------- Lücken, Linien, Karos ----------
% Einheitliche Maße für alle Blätter der Sequenz:
%  \linien{n}       n Schreiblinien, Abstand 8 mm, erste Linie 8 mm unter dem Text
%  \kariert{n}      Karoraster 5 mm, n Kästchen hoch, nur ganze Kästchen, mittig
%  \karobox{b}{h}   Karoraster b x h, auf ganze Kästchen gerundet
%  \luecke[l]       Lücke im Fließtext (Standard 4,5 cm, Fachbegriffe mind. 3,5 cm)
%  (jeweils mit Lösungsvariante \lsg…, siehe Kopf der Datei)
%  \restzeile       Lücke bis zum Zeilenende
%  \lueckentext     am Anfang eines Absatzes/Kastens mit Lücken: größerer
%                   Zeilenabstand und Flattersatz, damit man in die Lücken schreiben kann
%  \linienfeld{n}   umrahmtes Feld mit n Schreiblinien (z. B. für Java-Code)
%  \schreibzeile    Strut für Tabellenzellen, in die geschrieben wird (8 mm hoch)
\newlength{\zeilenabstand}\setlength{\zeilenabstand}{8mm}
\newlength{\karo}\setlength{\karo}{5mm}
\tikzset{karo/.style={draw=gray!50,line width=0.3pt},schreiblinie/.style={draw=black!65,line width=0.35pt}}
\newlength{\lsgbreite}
% Lücke: \luecke[Breite] / mit Lösung \lsg[Breite]{Text} (zu langer Text wird verkleinert)
\newcommand{\lsg}[2][4.5cm]{\leavevmode\ifloesung\settowidth{\lsgbreite}{\,#2\,}%
  \ifdim\lsgbreite>#1\relax\rlap{\raisebox{0.8pt}{\resizebox{#1}{!}{\lsgfarbe\,#2\,}}}%
  \else\rlap{\raisebox{0.8pt}{\lsgfarbe\,#2}}\fi\fi\rule[-0.6ex]{#1}{0.4pt}\xspace}
\newcommand{\luecke}[1][4.5cm]{\lsg[#1]{}}
% Lücke bis Zeilenende: \restzeile / \lsgrest{Text}
\newcommand{\lsgrest}[1]{\leavevmode\unskip\hspace{1.5mm}\ifloesung\rlap{\raisebox{0.8pt}{\lsgfarbe\,#1}}\fi\leaders\hbox{\rule[-0.6ex]{0.5pt}{0.4pt}}\hfill\null\par}
\newcommand{\restzeile}{\lsgrest{}}
% nur in der Lösung sichtbar (Platz bleibt frei): \lsgtext{…}; Alternativen: \lsgoder{Lösung}{Blatt}
\newcommand{\lsgtext}[1]{\ifloesung\leavevmode{\lsgfarbe#1}\fi}
\newcommand{\lsgoder}[2]{\ifloesung\leavevmode#1\else#2\fi}
\newcommand{\lueckentext}{\linespread{1.5}\selectfont\raggedright}
\newcommand{\schreibzeile}{\rule[-2.8mm]{0pt}{8mm}}
% Schritt-Nummern wie im Code und auf der Website: \nr{1} = ①
\newcommand{\nr}[1]{\ding{\numexpr171+#1\relax}}
% Schreiblinien: \linien{n} / \lsglinien{n}{Text} (Text liegt auf den Linien)
\newcommand{\lsglinien}[2]{\par\vspace{0.5mm}\noindent\begin{tikzpicture}
  \path (0,0) (\linewidth-0.4pt,-#1\zeilenabstand-1.5mm);
  \foreach \i in {1,...,#1}{\draw[schreiblinie] (0.2pt,-\i\zeilenabstand) -- (\linewidth-0.2pt,-\i\zeilenabstand);}
  \ifloesung\node[overlay,anchor=north west,inner sep=0pt] at (1mm,-\zeilenabstand+1pt+5mm) {\parbox[t]{\linewidth-2mm}{\lsgfarbe\raggedright\setlength{\baselineskip}{\zeilenabstand}\rule{0pt}{5mm}#2}};\fi
  \end{tikzpicture}\par}
\newcommand{\linien}[1]{\lsglinien{#1}{}}
% Codefeld mit Linien: \linienfeld{n} / \lsgfeld{n}{Code, Zeilen mit \\ getrennt}
\newcommand{\lsgfeld}[2]{\noindent\begin{tikzpicture}
  \draw[line width=0.4pt] (0.2pt,0) rectangle (\linewidth-0.2pt,-#1\zeilenabstand-3mm);
  \foreach \i in {1,...,#1}{\draw[schreiblinie] (2mm,-\i\zeilenabstand) -- (\linewidth-2mm,-\i\zeilenabstand);}
  \ifloesung\node[overlay,anchor=north west,inner sep=0pt] at (2.5mm,-\zeilenabstand+1pt+5mm) {\parbox[t]{\linewidth-5mm}{\lsgfarbe\ttfamily\small\setlength{\baselineskip}{\zeilenabstand}\rule{0pt}{5mm}#2}};\fi
  \end{tikzpicture}}
\newcommand{\linienfeld}[1]{\lsgfeld{#1}{}}
% Karoraster; Lösungszeichnung als TikZ-Code, Ursprung oben links, y nach unten negativ:
% \kariert{n} / \lsgkariert{n}{TikZ},  \karobox{b}{h} / \lsgkarobox{b}{h}{TikZ}
\newcommand{\karoraster}[3]{\begin{tikzpicture}\draw[karo] (0,0) grid[step=\karo] (#1\karo,#2\karo);
  \ifloesung\begin{scope}[overlay,shift={(0,#2\karo)},red!75!black,text=red!75!black,line width=0.8pt]#3\end{scope}\fi\end{tikzpicture}}
\newcommand{\lsgkariert}[2]{\par\vspace{1.5mm}\noindent\pgfmathtruncatemacro{\karoX}{floor(\linewidth/\karo)}%
  \makebox[\linewidth]{\karoraster{\karoX}{#1}{#2}}\par\vspace{1mm}}
\newcommand{\kariert}[1]{\lsgkariert{#1}{}}
\newcommand{\lsgkarobox}[3]{\pgfmathtruncatemacro{\karoX}{round((#1)/\karo)}\pgfmathtruncatemacro{\karoY}{round((#2)/\karo)}\karoraster{\karoX}{\karoY}{#3}}
\newcommand{\karobox}[2]{\lsgkarobox{#1}{#2}{}}
% Lösungs-Klassenkarte (für Zeichnungen im Karo): \lsgkarte{Name}{Attribute}{Methoden}
\newcommand{\lsgkarte}[3]{{\lsgfarbe\setlength{\arrayrulewidth}{0.6pt}\arrayrulecolor{red!75!black}\begin{tabular}{|l|}\hline\textbf{#1}\\\hline #2\\\hline #3\\\hline\end{tabular}}}
% Lösungs-Objektkarte (in TikZ als Knoteninhalt): \node[objekt] {\lsgobjekt{name : Klasse}{attr = wert\\…}};
\newcommand{\lsgobjekt}[2]{{\lsgfarbe\small\arrayrulecolor{red!75!black}\begin{tabular}{@{}l@{}}\textbf{#1}\\\hline #2\end{tabular}}}
\tikzset{objekt/.style={draw,rounded corners=3mm,inner sep=2mm,anchor=north west}}
\newcommand{\klassenkarte}[3]{\begin{tabular}{|l|}\hline\textbf{#1}\\\hline #2\\\hline #3\\\hline\end{tabular}}
\newcommand{\code}[1]{\texttt{#1}}

% ---------- Java-Listings ----------
\lstset{language=Java,basicstyle=\ttfamily\small,keywordstyle=\bfseries,commentstyle=\color{gray!80!black},stringstyle=\color{black},showstringspaces=false,columns=flexible,keepspaces=true,tabsize=3,frame=single,framerule=0.4pt,rulecolor=\color{black},xleftmargin=2mm,framexleftmargin=1mm,aboveskip=2mm,belowskip=2mm,breaklines=true,
 literate={ä}{{\"a}}1 {ö}{{\"o}}1 {ü}{{\"u}}1 {Ä}{{\"A}}1 {Ö}{{\"O}}1 {Ü}{{\"U}}1 {ß}{{\ss}}1 {–}{{--}}1 {„}{{\glqq}}1 {“}{{\grqq}}1}
\lstnewenvironment{java}[1][]{\lstset{#1}}{}
% Lückencode: größerer Zeilenabstand, Lücken mit (*\luecke[3cm]*) im Code
\lstnewenvironment{javaluecke}[1][]{\lstset{basicstyle=\linespread{1.5}\small\ttfamily,escapeinside={(*}{*)},#1}}{}

\newcolumntype{C}{>{\centering\arraybackslash}X}
\newcommand{\kopfzelle}[1]{\cellcolor{blue!8}\textbf{#1}}
\begin{document}
\abtitel{4.1}{Zahlensysteme}

\textbf{Das Dezimalsystem}\par
Unser gängiges Zahlensystem zur Basis 10 (griech. \textit{déka}, lat. \textit{decem}: zehn) kennt die zehn Ziffern 0 bis 9. Jede Stelle hat einen festen \textbf{Stellenwert}, eine Potenz von 10.\\[1mm]
Beispiel:\quad $35\,426 = 3\cdot 10\,000 + 5\cdot 1000 + 4\cdot 100 + 2\cdot 10 + 6\cdot 1$\\
\hspace*{2.05cm}$= 3\cdot 10^4 + 5\cdot 10^3 + 4\cdot 10^2 + 2\cdot 10^1 + 6\cdot 10^0$\par\vspace{1mm}
{\renewcommand{\arraystretch}{1.25}
\begin{tabularx}{\linewidth}{|>{\raggedright\arraybackslash}p{2.6cm}|C|C|C|C|C|}\hline
\kopfzelle{Stellenwert} & {\small Zehntausender ($10^4$)} & {\small Tausender ($10^3$)} & {\small Hunderter ($10^2$)} & {\small Zehner ($10^1$)} & {\small Einer ($10^0$)}\\\hline
\kopfzelle{Ziffer} & 3 & 5 & 4 & 2 & 6\\\hline
\end{tabularx}}

\medskip
\textbf{Das Binärsystem (zur Basis \lsg[1.5cm]{2})}\par
{\lueckentext Das Binärsystem kennt nur die Ziffern \lsg[1.5cm]{0} und \lsg[1.5cm]{1}. Eine solche Ziffer heißt \textbf{Bit} (\textit{binary digit}). Wie im Dezimalsystem hat jede Stelle einen festen \textbf{Stellenwert}. Die Stellenwerte werden \textbf{von rechts nach links} größer.\par}
\vspace{1mm}
{\renewcommand{\arraystretch}{1.1}
\begin{tabularx}{\linewidth}{|>{\raggedright\arraybackslash}p{2.6cm}|C|C|C|C|C|}\hline
\kopfzelle{Binärzahl} & \Large 1 & \Large 1 & \Large 1 & \Large 0 & \Large 1\\\hline
\kopfzelle{Stellenwert (Potenz)} & \schreibzeile\lsgtext{$2^4 = 16$} & \lsgtext{$2^3 = 8$} & \lsgtext{$2^2 = 4$} & \lsgtext{$2^1 = 2$} & \lsgtext{$2^0 = 1$}\\\hline
\kopfzelle{Produkt} & \schreibzeile\lsgtext{$1\cdot 16 = 16$} & \lsgtext{$1\cdot 8 = 8$} & \lsgtext{$1\cdot 4 = 4$} & \lsgtext{$0\cdot 2 = 0$} & \lsgtext{$1\cdot 1 = 1$}\\\hline
\kopfzelle{Dezimalzahl} & \multicolumn{5}{c|}{\schreibzeile\lsgtext{$16 + 8 + 4 + 0 + 1 = 29$}}\\\hline
\end{tabularx}}

\begin{merke}[Bit und Byte]
{\lueckentext Mit $n$ Bit lassen sich \lsg[2.5cm]{$2^n$} verschiedene Zahlen darstellen, nämlich 0 bis $2^n - 1$. 8 Bit fasst man zu \textbf{einem Byte} zusammen. Ein Byte kann also \lsg[2cm]{256} verschiedene Werte speichern, von 0 bis \lsg[1.5cm]{255}.\par}
\end{merke}

\begin{minipage}[t]{0.5\linewidth}
\textbf{Das Hexadezimalsystem (zur Basis \lsg[1.5cm]{16})}\par
Hier braucht man 16 Ziffern. Nach der 9 geht es mit Buchstaben weiter:\par\vspace{1.5mm}
{\small\setlength{\tabcolsep}{3pt}\renewcommand{\arraystretch}{1.15}
\begin{tabular}{|l|c|c|c|c|c|c|c|}\hline
\kopfzelle{Ziffer} & 0–9 & A & B & C & D & E & F\\\hline
\kopfzelle{Wert} & 0–9 & \schreibzeile\lsgtext{10} & \lsgtext{11} & \lsgtext{12} & \lsgtext{13} & \lsgtext{14} & \lsgtext{15}\\\hline
\end{tabular}}
\end{minipage}\hfill
\begin{minipage}[t]{0.47\linewidth}
\vspace{0pt}
{\renewcommand{\arraystretch}{1.1}\small
\begin{tabularx}{\linewidth}{|>{\raggedright\arraybackslash}p{2.35cm}|C|C|C|}\hline
\kopfzelle{Hexadezimal\-zahl} & \multicolumn{3}{c|}{\schreibzeile\Large\lsgtext{2\qquad B\qquad C}}\\\hline
\kopfzelle{Stellenwert (Potenz)} & \schreibzeile\lsgtext{$16^2{=}256$} & \lsgtext{$16^1{=}16$} & \lsgtext{$16^0{=}1$}\\\hline
\kopfzelle{Produkt} & \schreibzeile\lsgtext{$2\cdot 256$} & \lsgtext{$11\cdot 16$} & \lsgtext{$12\cdot 1$}\\\hline
\kopfzelle{Dezimalzahl} & \multicolumn{3}{c|}{\schreibzeile\normalsize 700}\\\hline
\end{tabularx}}
\end{minipage}

\medskip
\begin{minipage}[t]{0.45\linewidth}
\textbf{Umrechnung vom Dezimalsystem in ein anderes System} (Zahl $x$)\par\vspace{0.5mm}
\begin{enumerate}[label=\nr{\arabic*},leftmargin=6mm]
\item Finde den größten Stellenwert $S$ des neuen Systems, der nicht größer als $x$ ist.
\item Dividiere $x$ durch $S$ mit Rest $R$. Das ganzzahlige Ergebnis ist die Ziffer an dieser Stelle.
\item Mache mit dem Rest $R$ beim nächstkleineren Stellenwert weiter, bis die Stelle mit $S = 1$ erledigt ist.
\end{enumerate}
\end{minipage}\hfill
\begin{minipage}[t]{0.52\linewidth}
\vspace{0pt}
{\renewcommand{\arraystretch}{1.75}
\begin{tabular}{@{}l@{\ \,}l@{\ \,}l@{}}
\textbf{1. Durchgang:} & $S = 256$ & $700 : 256 = \mathbf{2}$ R $188$\\
\textbf{2. Durchgang:} & $S = $ \lsg[1.2cm]{16} & $188 : $ \lsg[3.2cm]{$16 = 11$ (B) R 12}\\
\textbf{3. Durchgang:} & $S = $ \lsg[1.2cm]{1} & \lsg[3.9cm]{$12 : 1 = 12$ (C) R 0}\\
\textbf{Ergebnis:} & \multicolumn{2}{l}{$700 = $ \lsg[2.5cm]{2BC}\textsubscript{16}}
\end{tabular}}
\end{minipage}

\newpage
\begin{aufgabe}[Bit-Schalter und Trainer]{1}
{\lueckentext\textbf{a)} Stelle mit dem \textbf{Bit-Schalter} auf der Website die Zahlen 5, 18 und 31 ein. Die größte Zahl, die man mit 5 Bit darstellen kann, ist \lsg[1.5cm]{31}, mit 8 Bit \lsg[1.5cm]{255}.\\
\textbf{b)} Bearbeite im \textbf{Zahlensystem-Trainer} (Level „Einfach“) mindestens 10 Aufgaben.\par}
\end{aufgabe}

\begin{aufgabe}[Umrechnen]{2}
Wandle die Zahlen \textbf{1337} und \textbf{4095} ins Hexadezimalsystem \textbf{und} ins Binärsystem um. Rechne mit dem Verfahren von Seite 1.
\lsgkariert{11}{\begin{scope}[font=\small,anchor=west]
\node at (2mm,-2.5mm) {1337 → Hex:\quad $1337 : 256 = 5$ R $57$;\quad $57 : 16 = 3$ R $9$;\quad $9 : 1 = 9$ \quad ⇒\quad 539\textsubscript{16}};
\node at (2mm,-7.5mm) {1337 → Binär: \ $S = 1024$: 1 R 313;\ 512: 0;\ 256: 1 R 57;\ 128: 0;\ 64: 0;\ 32: 1 R 25;};
\node at (2mm,-12.5mm) {\hspace{2.2cm}16: 1 R 9;\ 8: 1 R 1;\ 4: 0;\ 2: 0;\ 1: 1 \quad ⇒\quad 101\,0011\,1001\textsubscript{2}};
\node at (2mm,-22.5mm) {4095 → Hex:\quad $4095 : 256 = 15$ (F) R $255$;\quad $255 : 16 = 15$ (F) R $15$;\quad $15 : 1 = 15$ (F) \quad ⇒\quad FFF\textsubscript{16}};
\node at (2mm,-27.5mm) {4095 → Binär: \ $S = 2048$: 1 R 2047;\ 1024: 1 R 1023;\ 512: 1 R 511;\ \ldots\ jede Stelle 1};
\node at (2mm,-32.5mm) {\hspace{2.2cm}⇒\quad 1111\,1111\,1111\textsubscript{2}\quad (12 Bit; $4095 = 2^{12} - 1$)};
\end{scope}}
{\lueckentext $1337 = $ \lsg[3cm]{539}\textsubscript{16} $=$ \lsg[6cm]{101\,0011\,1001}\textsubscript{2}\\
$4095 = $ \lsg[3cm]{FFF}\textsubscript{16} $=$ \lsg[6cm]{1111\,1111\,1111}\textsubscript{2}\par}
\textit{Zum Nachdenken (mündlich):} Was fällt dir bei 4095 auf? Welche Zahl kommt danach?
\end{aufgabe}

\begin{aufgabe}[Die Abkürzung: Binär $\leftrightarrow$ Hexadezimal]{3}
Eine Hexadezimalziffer steht für genau 4 Bit, denn $16 = 2^4$. Teile eine Binärzahl deshalb \textbf{von rechts} in Vierergruppen und übersetze jede Gruppe einzeln.\\[1mm]
Beispiel:\quad \zs{1101\,0110}{2}\ $\to$\ 1101 = 13 = D,\ \ 0110 = 6\ $\to$\ \zs{D6}{16}\par\vspace{1mm}
{\lueckentext
\textbf{a)} \zs{1111\,0000}{2} = \lsg[1.6cm]{F0}\textsubscript{16}\qquad
\textbf{b)} \zs{1010\,0111}{2} = \lsg[1.6cm]{A7}\textsubscript{16}\qquad
\textbf{c)} \zs{7E}{16} = \lsg[3.2cm]{0111\,1110}\textsubscript{2}\\
\textbf{d)} Kontrolliere so dein Ergebnis aus Aufgabe 2 für 1337.\par}
\end{aufgabe}

\begin{aufgabe}[The Mystery of Cryptocastle]{4}
Öffne das Spiel \textit{The Mystery of Cryptocastle} (Link auf der Website) und bearbeite die Aufgaben in \textbf{„2 Zum alten Burgturm“}. Notiere hier deine Rechnungen und Codes:
\linien{3}
\end{aufgabe}

\vfill
\web{Station 4.1 Zahlensysteme: Bit-Schalter, Umrechnungs-Animation, Trainer}
\end{document}
;
%            füllt alle Lücken rot aus; siehe \lsg, \lsgrest, \lsglinien,
%            \lsgfeld, \lsgkariert, \lsgkarobox, \lsgtext, \lsgoder)
%  Lösungseinträge belegen nie mehr Platz als die Lücke im Blatt.
% =====================================================================
\usepackage[a4paper,left=18mm,right=18mm,top=26mm,bottom=17mm,headheight=15mm,headsep=4mm,footskip=8mm]{geometry}
\usepackage{iftex}
\ifPDFTeX
  \usepackage[T1]{fontenc}
  \usepackage[utf8]{inputenc}
  \usepackage{helvet}
  \usepackage[scaled=0.86]{beramono}
  \renewcommand{\familydefault}{\sfdefault}
\else
  \usepackage{fontspec}
  \setmainfont{texgyreheros}[Extension=.otf,UprightFont=*-regular,BoldFont=*-bold,ItalicFont=*-italic,BoldItalicFont=*-bolditalic]
  \setmonofont{DejaVuSansMono}[Extension=.ttf,UprightFont=*,BoldFont=*-Bold,ItalicFont=*-Oblique,BoldItalicFont=*-BoldOblique,Scale=0.86]
\fi
\usepackage[ngerman,shorthands=off]{babel}
\usepackage{xcolor,graphicx,tabularx,array,enumitem,fancyhdr,tikz,listings,multicol,calc,etoolbox,ragged2e,colortbl,xspace,pifont}
\usepackage[most]{tcolorbox}
\usetikzlibrary{arrows.meta,positioning,calc,decorations.pathreplacing,shapes.geometric,shapes.symbols,fit}
\usepackage[hidelinks]{hyperref}
\setlength{\parindent}{0pt}
\setlength{\parskip}{3pt}
\setlist{nosep,leftmargin=*}

% ---------- Lösungsschalter ----------
\newif\ifloesung
\ifdefined\mitloesung\loesungtrue\fi
\newcommand{\lsgfarbe}{\color{red!75!black}}

% ---------- Kopf- und Fußzeile (einheitliche Form) ----------
\newcommand{\lehrkraft}{Hau}
\newcommand{\themenbereich}{4. Codierung}
\newcommand{\abnummer}{}
\pagestyle{fancy}\fancyhf{}
\renewcommand{\headrulewidth}{0.4pt}
\fancyhead[L]{\small Klasse 11\\Datum:}
\fancyhead[C]{\small Themenbereich:\\\themenbereich}
\fancyhead[R]{\small \lehrkraft\ifloesung\\{\lsgfarbe\bfseries Lösung}\fi}
\fancyfoot[C]{\scriptsize edu-mrh.de\,·\,Informatik 11\,·\,Blatt \abnummer\ifloesung\ (Lösung)\fi\,·\,Seite \thepage}
\newcommand{\abtitel}[2]{\renewcommand{\abnummer}{#1}\begin{center}\Large\bfseries #1\quad\underline{#2}\end{center}\vspace{-1mm}}

% ---------- Boxen ----------
\newenvironment{aufgabe}[2][]{\begin{tcolorbox}[enhanced,breakable,colback=white,colframe=black,boxrule=0.7pt,arc=3mm,left=2mm,right=2mm,top=1.5mm,bottom=1.5mm,before skip=2.5mm,after skip=2.5mm]\textbf{Aufgabe #2}\ifstrempty{#1}{}{ \textit{#1}}\par\vspace{0.6mm}}{\end{tcolorbox}}
\newtcolorbox{merke}[1][Merke]{enhanced,breakable,colback=green!5,colframe=green!40!black,boxrule=0.8pt,arc=2mm,left=2mm,right=2mm,top=1.5mm,bottom=1.5mm,before skip=2.5mm,after skip=2.5mm,title={#1},fonttitle=\bfseries,coltitle=white}
\newtcolorbox{hinweis}{enhanced,breakable,colback=yellow!12,colframe=orange!70!black,boxrule=0.6pt,arc=2mm,left=2mm,right=2mm,top=1mm,bottom=1mm,before skip=2mm,after skip=2mm}
\newcommand{\web}[1]{{\small\color{gray}\textit{Website: #1}}}
% Kasten für optionale Aufgaben
\newtcolorbox{optional}[1][Optional]{enhanced,breakable,colback=blue!3,colframe=blue!45!black,boxrule=0.6pt,arc=2mm,left=2mm,right=2mm,top=1.5mm,bottom=1.5mm,before skip=2.5mm,after skip=2.5mm,title={#1},fonttitle=\bfseries\small,coltitle=white}
% Zahl mit Basis: \zs{1011}{2}
\newcommand{\zs}[2]{#1\textsubscript{#2}}

% ---------- Lücken, Linien, Karos ----------
% Einheitliche Maße für alle Blätter der Sequenz:
%  \linien{n}       n Schreiblinien, Abstand 8 mm, erste Linie 8 mm unter dem Text
%  \kariert{n}      Karoraster 5 mm, n Kästchen hoch, nur ganze Kästchen, mittig
%  \karobox{b}{h}   Karoraster b x h, auf ganze Kästchen gerundet
%  \luecke[l]       Lücke im Fließtext (Standard 4,5 cm, Fachbegriffe mind. 3,5 cm)
%  (jeweils mit Lösungsvariante \lsg…, siehe Kopf der Datei)
%  \restzeile       Lücke bis zum Zeilenende
%  \lueckentext     am Anfang eines Absatzes/Kastens mit Lücken: größerer
%                   Zeilenabstand und Flattersatz, damit man in die Lücken schreiben kann
%  \linienfeld{n}   umrahmtes Feld mit n Schreiblinien (z. B. für Java-Code)
%  \schreibzeile    Strut für Tabellenzellen, in die geschrieben wird (8 mm hoch)
\newlength{\zeilenabstand}\setlength{\zeilenabstand}{8mm}
\newlength{\karo}\setlength{\karo}{5mm}
\tikzset{karo/.style={draw=gray!50,line width=0.3pt},schreiblinie/.style={draw=black!65,line width=0.35pt}}
\newlength{\lsgbreite}
% Lücke: \luecke[Breite] / mit Lösung \lsg[Breite]{Text} (zu langer Text wird verkleinert)
\newcommand{\lsg}[2][4.5cm]{\leavevmode\ifloesung\settowidth{\lsgbreite}{\,#2\,}%
  \ifdim\lsgbreite>#1\relax\rlap{\raisebox{0.8pt}{\resizebox{#1}{!}{\lsgfarbe\,#2\,}}}%
  \else\rlap{\raisebox{0.8pt}{\lsgfarbe\,#2}}\fi\fi\rule[-0.6ex]{#1}{0.4pt}\xspace}
\newcommand{\luecke}[1][4.5cm]{\lsg[#1]{}}
% Lücke bis Zeilenende: \restzeile / \lsgrest{Text}
\newcommand{\lsgrest}[1]{\leavevmode\unskip\hspace{1.5mm}\ifloesung\rlap{\raisebox{0.8pt}{\lsgfarbe\,#1}}\fi\leaders\hbox{\rule[-0.6ex]{0.5pt}{0.4pt}}\hfill\null\par}
\newcommand{\restzeile}{\lsgrest{}}
% nur in der Lösung sichtbar (Platz bleibt frei): \lsgtext{…}; Alternativen: \lsgoder{Lösung}{Blatt}
\newcommand{\lsgtext}[1]{\ifloesung\leavevmode{\lsgfarbe#1}\fi}
\newcommand{\lsgoder}[2]{\ifloesung\leavevmode#1\else#2\fi}
\newcommand{\lueckentext}{\linespread{1.5}\selectfont\raggedright}
\newcommand{\schreibzeile}{\rule[-2.8mm]{0pt}{8mm}}
% Schritt-Nummern wie im Code und auf der Website: \nr{1} = ①
\newcommand{\nr}[1]{\ding{\numexpr171+#1\relax}}
% Schreiblinien: \linien{n} / \lsglinien{n}{Text} (Text liegt auf den Linien)
\newcommand{\lsglinien}[2]{\par\vspace{0.5mm}\noindent\begin{tikzpicture}
  \path (0,0) (\linewidth-0.4pt,-#1\zeilenabstand-1.5mm);
  \foreach \i in {1,...,#1}{\draw[schreiblinie] (0.2pt,-\i\zeilenabstand) -- (\linewidth-0.2pt,-\i\zeilenabstand);}
  \ifloesung\node[overlay,anchor=north west,inner sep=0pt] at (1mm,-\zeilenabstand+1pt+5mm) {\parbox[t]{\linewidth-2mm}{\lsgfarbe\raggedright\setlength{\baselineskip}{\zeilenabstand}\rule{0pt}{5mm}#2}};\fi
  \end{tikzpicture}\par}
\newcommand{\linien}[1]{\lsglinien{#1}{}}
% Codefeld mit Linien: \linienfeld{n} / \lsgfeld{n}{Code, Zeilen mit \\ getrennt}
\newcommand{\lsgfeld}[2]{\noindent\begin{tikzpicture}
  \draw[line width=0.4pt] (0.2pt,0) rectangle (\linewidth-0.2pt,-#1\zeilenabstand-3mm);
  \foreach \i in {1,...,#1}{\draw[schreiblinie] (2mm,-\i\zeilenabstand) -- (\linewidth-2mm,-\i\zeilenabstand);}
  \ifloesung\node[overlay,anchor=north west,inner sep=0pt] at (2.5mm,-\zeilenabstand+1pt+5mm) {\parbox[t]{\linewidth-5mm}{\lsgfarbe\ttfamily\small\setlength{\baselineskip}{\zeilenabstand}\rule{0pt}{5mm}#2}};\fi
  \end{tikzpicture}}
\newcommand{\linienfeld}[1]{\lsgfeld{#1}{}}
% Karoraster; Lösungszeichnung als TikZ-Code, Ursprung oben links, y nach unten negativ:
% \kariert{n} / \lsgkariert{n}{TikZ},  \karobox{b}{h} / \lsgkarobox{b}{h}{TikZ}
\newcommand{\karoraster}[3]{\begin{tikzpicture}\draw[karo] (0,0) grid[step=\karo] (#1\karo,#2\karo);
  \ifloesung\begin{scope}[overlay,shift={(0,#2\karo)},red!75!black,text=red!75!black,line width=0.8pt]#3\end{scope}\fi\end{tikzpicture}}
\newcommand{\lsgkariert}[2]{\par\vspace{1.5mm}\noindent\pgfmathtruncatemacro{\karoX}{floor(\linewidth/\karo)}%
  \makebox[\linewidth]{\karoraster{\karoX}{#1}{#2}}\par\vspace{1mm}}
\newcommand{\kariert}[1]{\lsgkariert{#1}{}}
\newcommand{\lsgkarobox}[3]{\pgfmathtruncatemacro{\karoX}{round((#1)/\karo)}\pgfmathtruncatemacro{\karoY}{round((#2)/\karo)}\karoraster{\karoX}{\karoY}{#3}}
\newcommand{\karobox}[2]{\lsgkarobox{#1}{#2}{}}
% Lösungs-Klassenkarte (für Zeichnungen im Karo): \lsgkarte{Name}{Attribute}{Methoden}
\newcommand{\lsgkarte}[3]{{\lsgfarbe\setlength{\arrayrulewidth}{0.6pt}\arrayrulecolor{red!75!black}\begin{tabular}{|l|}\hline\textbf{#1}\\\hline #2\\\hline #3\\\hline\end{tabular}}}
% Lösungs-Objektkarte (in TikZ als Knoteninhalt): \node[objekt] {\lsgobjekt{name : Klasse}{attr = wert\\…}};
\newcommand{\lsgobjekt}[2]{{\lsgfarbe\small\arrayrulecolor{red!75!black}\begin{tabular}{@{}l@{}}\textbf{#1}\\\hline #2\end{tabular}}}
\tikzset{objekt/.style={draw,rounded corners=3mm,inner sep=2mm,anchor=north west}}
\newcommand{\klassenkarte}[3]{\begin{tabular}{|l|}\hline\textbf{#1}\\\hline #2\\\hline #3\\\hline\end{tabular}}
\newcommand{\code}[1]{\texttt{#1}}

% ---------- Java-Listings ----------
\lstset{language=Java,basicstyle=\ttfamily\small,keywordstyle=\bfseries,commentstyle=\color{gray!80!black},stringstyle=\color{black},showstringspaces=false,columns=flexible,keepspaces=true,tabsize=3,frame=single,framerule=0.4pt,rulecolor=\color{black},xleftmargin=2mm,framexleftmargin=1mm,aboveskip=2mm,belowskip=2mm,breaklines=true,
 literate={ä}{{\"a}}1 {ö}{{\"o}}1 {ü}{{\"u}}1 {Ä}{{\"A}}1 {Ö}{{\"O}}1 {Ü}{{\"U}}1 {ß}{{\ss}}1 {–}{{--}}1 {„}{{\glqq}}1 {“}{{\grqq}}1}
\lstnewenvironment{java}[1][]{\lstset{#1}}{}
% Lückencode: größerer Zeilenabstand, Lücken mit (*\luecke[3cm]*) im Code
\lstnewenvironment{javaluecke}[1][]{\lstset{basicstyle=\linespread{1.5}\small\ttfamily,escapeinside={(*}{*)},#1}}{}

\newcolumntype{C}{>{\centering\arraybackslash}X}
\newcommand{\kopfzelle}[1]{\cellcolor{blue!8}\textbf{#1}}
\begin{document}
\abtitel{4.1}{Zahlensysteme}

\textbf{Das Dezimalsystem}\par
Unser gängiges Zahlensystem zur Basis 10 (griech. \textit{déka}, lat. \textit{decem}: zehn) kennt die zehn Ziffern 0 bis 9. Jede Stelle hat einen festen \textbf{Stellenwert}, eine Potenz von 10.\\[1mm]
Beispiel:\quad $35\,426 = 3\cdot 10\,000 + 5\cdot 1000 + 4\cdot 100 + 2\cdot 10 + 6\cdot 1$\\
\hspace*{2.05cm}$= 3\cdot 10^4 + 5\cdot 10^3 + 4\cdot 10^2 + 2\cdot 10^1 + 6\cdot 10^0$\par\vspace{1mm}
{\renewcommand{\arraystretch}{1.25}
\begin{tabularx}{\linewidth}{|>{\raggedright\arraybackslash}p{2.6cm}|C|C|C|C|C|}\hline
\kopfzelle{Stellenwert} & {\small Zehntausender ($10^4$)} & {\small Tausender ($10^3$)} & {\small Hunderter ($10^2$)} & {\small Zehner ($10^1$)} & {\small Einer ($10^0$)}\\\hline
\kopfzelle{Ziffer} & 3 & 5 & 4 & 2 & 6\\\hline
\end{tabularx}}

\medskip
\textbf{Das Binärsystem (zur Basis \lsg[1.5cm]{2})}\par
{\lueckentext Das Binärsystem kennt nur die Ziffern \lsg[1.5cm]{0} und \lsg[1.5cm]{1}. Eine solche Ziffer heißt \textbf{Bit} (\textit{binary digit}). Wie im Dezimalsystem hat jede Stelle einen festen \textbf{Stellenwert}. Die Stellenwerte werden \textbf{von rechts nach links} größer.\par}
\vspace{1mm}
{\renewcommand{\arraystretch}{1.1}
\begin{tabularx}{\linewidth}{|>{\raggedright\arraybackslash}p{2.6cm}|C|C|C|C|C|}\hline
\kopfzelle{Binärzahl} & \Large 1 & \Large 1 & \Large 1 & \Large 0 & \Large 1\\\hline
\kopfzelle{Stellenwert (Potenz)} & \schreibzeile\lsgtext{$2^4 = 16$} & \lsgtext{$2^3 = 8$} & \lsgtext{$2^2 = 4$} & \lsgtext{$2^1 = 2$} & \lsgtext{$2^0 = 1$}\\\hline
\kopfzelle{Produkt} & \schreibzeile\lsgtext{$1\cdot 16 = 16$} & \lsgtext{$1\cdot 8 = 8$} & \lsgtext{$1\cdot 4 = 4$} & \lsgtext{$0\cdot 2 = 0$} & \lsgtext{$1\cdot 1 = 1$}\\\hline
\kopfzelle{Dezimalzahl} & \multicolumn{5}{c|}{\schreibzeile\lsgtext{$16 + 8 + 4 + 0 + 1 = 29$}}\\\hline
\end{tabularx}}

\begin{merke}[Bit und Byte]
{\lueckentext Mit $n$ Bit lassen sich \lsg[2.5cm]{$2^n$} verschiedene Zahlen darstellen, nämlich 0 bis $2^n - 1$. 8 Bit fasst man zu \textbf{einem Byte} zusammen. Ein Byte kann also \lsg[2cm]{256} verschiedene Werte speichern, von 0 bis \lsg[1.5cm]{255}.\par}
\end{merke}

\begin{minipage}[t]{0.5\linewidth}
\textbf{Das Hexadezimalsystem (zur Basis \lsg[1.5cm]{16})}\par
Hier braucht man 16 Ziffern. Nach der 9 geht es mit Buchstaben weiter:\par\vspace{1.5mm}
{\small\setlength{\tabcolsep}{3pt}\renewcommand{\arraystretch}{1.15}
\begin{tabular}{|l|c|c|c|c|c|c|c|}\hline
\kopfzelle{Ziffer} & 0–9 & A & B & C & D & E & F\\\hline
\kopfzelle{Wert} & 0–9 & \schreibzeile\lsgtext{10} & \lsgtext{11} & \lsgtext{12} & \lsgtext{13} & \lsgtext{14} & \lsgtext{15}\\\hline
\end{tabular}}
\end{minipage}\hfill
\begin{minipage}[t]{0.47\linewidth}
\vspace{0pt}
{\renewcommand{\arraystretch}{1.1}\small
\begin{tabularx}{\linewidth}{|>{\raggedright\arraybackslash}p{2.35cm}|C|C|C|}\hline
\kopfzelle{Hexadezimal\-zahl} & \multicolumn{3}{c|}{\schreibzeile\Large\lsgtext{2\qquad B\qquad C}}\\\hline
\kopfzelle{Stellenwert (Potenz)} & \schreibzeile\lsgtext{$16^2{=}256$} & \lsgtext{$16^1{=}16$} & \lsgtext{$16^0{=}1$}\\\hline
\kopfzelle{Produkt} & \schreibzeile\lsgtext{$2\cdot 256$} & \lsgtext{$11\cdot 16$} & \lsgtext{$12\cdot 1$}\\\hline
\kopfzelle{Dezimalzahl} & \multicolumn{3}{c|}{\schreibzeile\normalsize 700}\\\hline
\end{tabularx}}
\end{minipage}

\medskip
\begin{minipage}[t]{0.45\linewidth}
\textbf{Umrechnung vom Dezimalsystem in ein anderes System} (Zahl $x$)\par\vspace{0.5mm}
\begin{enumerate}[label=\nr{\arabic*},leftmargin=6mm]
\item Finde den größten Stellenwert $S$ des neuen Systems, der nicht größer als $x$ ist.
\item Dividiere $x$ durch $S$ mit Rest $R$. Das ganzzahlige Ergebnis ist die Ziffer an dieser Stelle.
\item Mache mit dem Rest $R$ beim nächstkleineren Stellenwert weiter, bis die Stelle mit $S = 1$ erledigt ist.
\end{enumerate}
\end{minipage}\hfill
\begin{minipage}[t]{0.52\linewidth}
\vspace{0pt}
{\renewcommand{\arraystretch}{1.75}
\begin{tabular}{@{}l@{\ \,}l@{\ \,}l@{}}
\textbf{1. Durchgang:} & $S = 256$ & $700 : 256 = \mathbf{2}$ R $188$\\
\textbf{2. Durchgang:} & $S = $ \lsg[1.2cm]{16} & $188 : $ \lsg[3.2cm]{$16 = 11$ (B) R 12}\\
\textbf{3. Durchgang:} & $S = $ \lsg[1.2cm]{1} & \lsg[3.9cm]{$12 : 1 = 12$ (C) R 0}\\
\textbf{Ergebnis:} & \multicolumn{2}{l}{$700 = $ \lsg[2.5cm]{2BC}\textsubscript{16}}
\end{tabular}}
\end{minipage}

\newpage
\begin{aufgabe}[Bit-Schalter und Trainer]{1}
{\lueckentext\textbf{a)} Stelle mit dem \textbf{Bit-Schalter} auf der Website die Zahlen 5, 18 und 31 ein. Die größte Zahl, die man mit 5 Bit darstellen kann, ist \lsg[1.5cm]{31}, mit 8 Bit \lsg[1.5cm]{255}.\\
\textbf{b)} Bearbeite im \textbf{Zahlensystem-Trainer} (Level „Einfach“) mindestens 10 Aufgaben.\par}
\end{aufgabe}

\begin{aufgabe}[Umrechnen]{2}
Wandle die Zahlen \textbf{1337} und \textbf{4095} ins Hexadezimalsystem \textbf{und} ins Binärsystem um. Rechne mit dem Verfahren von Seite 1.
\lsgkariert{11}{\begin{scope}[font=\small,anchor=west]
\node at (2mm,-2.5mm) {1337 → Hex:\quad $1337 : 256 = 5$ R $57$;\quad $57 : 16 = 3$ R $9$;\quad $9 : 1 = 9$ \quad ⇒\quad 539\textsubscript{16}};
\node at (2mm,-7.5mm) {1337 → Binär: \ $S = 1024$: 1 R 313;\ 512: 0;\ 256: 1 R 57;\ 128: 0;\ 64: 0;\ 32: 1 R 25;};
\node at (2mm,-12.5mm) {\hspace{2.2cm}16: 1 R 9;\ 8: 1 R 1;\ 4: 0;\ 2: 0;\ 1: 1 \quad ⇒\quad 101\,0011\,1001\textsubscript{2}};
\node at (2mm,-22.5mm) {4095 → Hex:\quad $4095 : 256 = 15$ (F) R $255$;\quad $255 : 16 = 15$ (F) R $15$;\quad $15 : 1 = 15$ (F) \quad ⇒\quad FFF\textsubscript{16}};
\node at (2mm,-27.5mm) {4095 → Binär: \ $S = 2048$: 1 R 2047;\ 1024: 1 R 1023;\ 512: 1 R 511;\ \ldots\ jede Stelle 1};
\node at (2mm,-32.5mm) {\hspace{2.2cm}⇒\quad 1111\,1111\,1111\textsubscript{2}\quad (12 Bit; $4095 = 2^{12} - 1$)};
\end{scope}}
{\lueckentext $1337 = $ \lsg[3cm]{539}\textsubscript{16} $=$ \lsg[6cm]{101\,0011\,1001}\textsubscript{2}\\
$4095 = $ \lsg[3cm]{FFF}\textsubscript{16} $=$ \lsg[6cm]{1111\,1111\,1111}\textsubscript{2}\par}
\textit{Zum Nachdenken (mündlich):} Was fällt dir bei 4095 auf? Welche Zahl kommt danach?
\end{aufgabe}

\begin{aufgabe}[Die Abkürzung: Binär $\leftrightarrow$ Hexadezimal]{3}
Eine Hexadezimalziffer steht für genau 4 Bit, denn $16 = 2^4$. Teile eine Binärzahl deshalb \textbf{von rechts} in Vierergruppen und übersetze jede Gruppe einzeln.\\[1mm]
Beispiel:\quad \zs{1101\,0110}{2}\ $\to$\ 1101 = 13 = D,\ \ 0110 = 6\ $\to$\ \zs{D6}{16}\par\vspace{1mm}
{\lueckentext
\textbf{a)} \zs{1111\,0000}{2} = \lsg[1.6cm]{F0}\textsubscript{16}\qquad
\textbf{b)} \zs{1010\,0111}{2} = \lsg[1.6cm]{A7}\textsubscript{16}\qquad
\textbf{c)} \zs{7E}{16} = \lsg[3.2cm]{0111\,1110}\textsubscript{2}\\
\textbf{d)} Kontrolliere so dein Ergebnis aus Aufgabe 2 für 1337.\par}
\end{aufgabe}

\begin{aufgabe}[The Mystery of Cryptocastle]{4}
Öffne das Spiel \textit{The Mystery of Cryptocastle} (Link auf der Website) und bearbeite die Aufgaben in \textbf{„2 Zum alten Burgturm“}. Notiere hier deine Rechnungen und Codes:
\linien{3}
\end{aufgabe}

\vfill
\web{Station 4.1 Zahlensysteme: Bit-Schalter, Umrechnungs-Animation, Trainer}
\end{document}
